\documentclass{article}

\usepackage[italian]{babel} %testi autogenerati italiano
\usepackage{minted}         %per codice vhdl bello
\usepackage{tikz}           %per disegno fsm
\usetikzlibrary{automata, positioning, arrows}
\usepackage{circuitikz}     %per disegno componente
\usepackage{graphicx}       %per importare immagini
\usepackage{geometry}       %per gestire margini e spostamenti
\geometry {
    top=25mm,
    bmargin=25mm,
}
\usepackage{array}          %per colonne di width fissata
\usepackage{subcaption}     %tabelle divise
\usepackage{hyperref}       %links
\hypersetup{
    colorlinks=true,
    linkcolor=blue,
}

\begin{document}

\setlength\parindent{0pt} %noindent automatico
\setlength\parskip{1em}

\begin{titlepage}
    \centering
    \hrule

    \vspace{0,5cm}
    {
        \normalsize Politecnico di Milano\\
        Dipartimento di Elettronica, Informazione e Bioingegneria
    }

    \vspace{4,1cm}
    {\Huge \textbf{Progetto di Reti Logiche\\
            2020/21}\\}

    \vspace{0,5cm}
    \large {Prof. Palermo Gianluca}

    \vspace{2,5cm}
    {
        \large
        \begin{tabular}{c c}
            Shalby Hazem Hesham Yousef & (Codice Persona: 10596243, Matricola: 910871) \\
            Perego Niccolò             & (Codice Persona: 10628782, Matricola: 895468) \\
        \end{tabular}

    }

    \vspace{3,8cm}

    \normalsize{1 Aprile 2021}
    \vspace{0,2cm}

    \centering\hspace{0,2cm}\includegraphics[scale=0.6]{logo.png}
    \vspace{0,5cm}
    \hrule

\end{titlepage}

\pagebreak

\section{Requisiti di progetto} %1

\subsection{Descrizione del problema} %1.1
Si vuole realizzare un componente in grado di svolgere una versione semplificata del processo di equalizzazione dell’istogramma di un’immagine, ossia di ricalibrare il contrasto di quest’ultima,
effettuando una ridistribuzione dei valori di intensità pixel per pixel. \par
Le immagini di cui è richiesta la manipolazione sono definite in scala di grigi a 256 livelli e hanno una dimensione massima di 128x128 pixel.
\vspace{0,3cm} %un po' di spazio

\subsection{Interfaccia del componente} %1.2
Il componente realizzato ha un’interfaccia così definita in liguaggio VHDL:

\begin{minted}{vhdl}
    ENTITY project_reti_logiche IS PORT
	(
		i_clk     : IN  std_logic;
		i_rst     : IN  std_logic; 
		i_start   : IN  std_logic; 
		i_data    : IN  std_logic_vector (7 DOWNTO 0);
		o_address : OUT std_logic_vector (15 DOWNTO 0);
		o_done    : OUT std_logic; 
		o_en      : OUT std_logic;
		o_we      : OUT std_logic;
		o_data    : OUT std_logic_vector (7 DOWNTO 0) 
	);
    END project_reti_logiche;
\end{minted}

\vspace{0,3cm} %un po' di spazio

In particolare:
\begin{itemize}
    \item   \texttt{i\_clk} è il segnale di CLOCK in ingresso generato dal TestBench;
    \item   \texttt{i\_rst} è il segnale di RESET che inizializza la macchina, predisponendola alla ricezione del segnale di START. Può essere anche asincrono;
    \item   \texttt{i\_start} è il segnale di START generato dal Test Bench;
    \item   \texttt{i\_data} è il segnale (vettore) che arriva dalla memoria in seguito a una richiesta di lettura;
    \item   \texttt{o\_address} è il segnale (vettore) di uscita che manda l’indirizzo alla memoria;
    \item   \texttt{o\_done} è il segnale di uscita che comunica la fine dell’elaborazione e il dato di uscita scritto in memoria;
    \item   \texttt{o\_en} è il segnale di ENABLE da mandare alla memoria per poter comunicare (sia in lettura che in scrittura);
    \item   \texttt{o\_we} è il segnale di WRITE ENABLE da mandare alla memoria per comunicare quale operazione si vuole svolgere su di essa. Può assumere valori 0 e 1, rispettivamente per lettura e scrittura;
    \item   \texttt{o\_data} è il segnale (vettore) di uscita dal componente verso la memoria.
\end{itemize}

\pagebreak

\begin{figure}[ht]

    \centering
    \begin{circuitikz}
        \ctikzset{multipoles/thickness=3.5}
        \ctikzset{multipoles/dipchip/width=4}
        %sinistra
        \draw   (0,0) node[dipchip,
            num pins=10, hide numbers, no topmark,
            external pins width=0](C){\large{\texttt{COMPONENT}}};
        \node   [right] at (C.bpin 1) {\texttt{i\_start}};
        \draw   (C.bpin 1) -- ++(-1,0) coordinate(extpin);
        \node   [right] at (C.bpin 2) {\texttt{i\_data}};
        \draw   (C.bpin 2) -- ++(-1,0) coordinate(extpin);
        \node   [right] at (C.bpin 4) {\texttt{i\_rst}};
        \draw   (C.bpin 4) -- ++(-1,0) coordinate(extpin);
        \draw   (C.bpin 5) ++(0,0.2) -- ++(0.2,-0.2)
        node[right]{\texttt{i\_clk}} -- ++(-0.2,-0.2);
        \draw   (C.bpin 5) -- ++(-1,0) coordinate(extpin);
        %destra
        \node   [left] at (C.bpin 6) {\texttt{o\_data}};
        \draw   (C.bpin 6) -- ++(1,0) coordinate(extpin);
        \node   [left] at (C.bpin 7) {\texttt{o\_done}};
        \draw   (C.bpin 7) -- ++(1,0) coordinate(extpin);
        \node   [left] at (C.bpin 8) {\texttt{o\_en}};
        \draw   (C.bpin 8) -- ++(1,0) coordinate(extpin);
        \node   [left] at (C.bpin 9) {\texttt{o\_we}};
        \draw   (C.bpin 9) -- ++(1,0) coordinate(extpin);
        \node   [left] at (C.bpin 10) {\texttt{o\_address}};
        \draw   (C.bpin 10) -- ++(1,0) coordinate(extpin);
    \end{circuitikz}

    \caption{Schema del componente realizzato.}
    \label{fig:component}
\end{figure}
\vspace{0,3cm}

\subsection{Descrizione della memoria e dell'interazione con il componente} %1.3
Il modulo implementato dovrà dialogare in lettura e scrittura con una RAM, indirizzata al byte.\par
In particolare, l’algoritmo di equalizzazione sarà applicato a immagini pre-salvate in memoria, la cui grandezza effettiva (in pixel) sarà specificata dal prodotto del contenuto tra le celle a indirizzo \texttt{0} e \texttt{1} della RAM, contenenti rispettivamente il numero di colonne \texttt{n\_col} e di righe \texttt{n\_row} dell’immagine, entrambi di dimensione 8 bit.\\
Nei byte successivi, dall’indirizzo \texttt{2} all’indirizzo \texttt{n\_col $\cdot$ n\_row $+$ 1}, sarà contenuta, pixel per pixel, sequenzialmente e in modo contiguo, l’immagine di cui è richiesta la trasformazione.\\
Infine, l’immagine ottenuta dal processo di equalizzazione svolto dal componente verrà salvata in memoria dall’indirizzo \texttt{n\_col $\cdot$ n\_row $+$ 2} all’indirizzo \texttt{2 $\cdot$ n\_col $\cdot$ n\_row $+$ 1}.
\vspace{0,3cm} %un po' di spazio

\begin{table}[h]
    \centering
    \def\arraystretch{1.3} %un po' di padding

    \caption{Schema generale del contenuto della memoria dopo l'elaborazione.}

    \begin{tabular}{ |m{4cm} | m{5cm} }
        \cline{1-1}
        \texttt{n\_col}                             & \texttt{addr. 0}                                     \\
        \cline{1-1}
        \texttt{n\_row}                             & \texttt{addr. 1}                                     \\
        \cline{1-1}
        \texttt{Primo pixel immagine da elaborare}  & \texttt{addr. 2}                                     \\
        \cline{1-1}
        \vdots                                      & \vdots                                               \\
        \cline{1-1}
        \texttt{Ultimo pixel immagine da elaborare} & \texttt{addr. n\_col $\cdot$ n\_row $+$ 1}           \\
        \cline{1-1}
        \texttt{Primo pixel immagine elaborata}     & \texttt{addr. n\_col $\cdot$ n\_row $+$ 2}           \\
        \cline{1-1}
        \vdots                                      & \vdots                                               \\
        \cline{1-1}
        \texttt{Ultimo pixel immagine elaborata}    & \texttt{addr. 2 $\cdot$ n\_col $\cdot$ n\_row $+$ 1} \\
        \cline{1-1}
    \end{tabular}

    \label{tab:schemamem}

\end{table}

\vspace{0.3cm}

\subsection{Esempio di funzionamento} %1.4
\label{sec:esempio}
Si riporta in seguito un'esempio di elaborazione compiuta su un'immagine di test.

\pagebreak

\begin{figure}[ht]
    \hspace{-10pt}
    \begin{minipage}[c]{.35\linewidth}
        \centering
        \begin{tabular}{c|c}
            addr. & data \\
            \hline \hline
            0     & 4    \\
            1     & 4    \\
            2     & 52   \\
            3     & 55   \\
            4     & 122  \\
            5     & 70   \\
            6     & 62   \\
            7     & 59   \\
            8     & 68   \\
            9     & 113  \\
            10    & 85   \\
            11    & 71   \\
            12    & 59   \\
            13    & 64   \\
            14    & 87   \\
            15    & 78   \\
            16    & 68   \\
        \end{tabular}
        \begin{tabular}{c|c}
            addr. & data \\
            \hline \hline
            17    & 69   \\
            18    & 0    \\
            19    & 12   \\
            20    & 255  \\
            21    & 72   \\
            22    & 40   \\
            23    & 28   \\
            24    & 64   \\
            25    & 244  \\
            26    & 132  \\
            27    & 76   \\
            28    & 28   \\
            29    & 48   \\
            30    & 140  \\
            31    & 104  \\
            32    & 64   \\
            33    & 68   \\
        \end{tabular}
        \subcaption{Contenuto della memoria}\label{esempio1}
    \end{minipage}%
    \begin{minipage}[c]{.3\linewidth}
        \centering
        \includegraphics[scale=0.25]{immSorg.jpg}
        \subcaption{Immagine sorgente}\label{esempio2}
    \end{minipage}\hspace{20pt}
    \begin{minipage}[c]{.3\linewidth}
        \centering
        \includegraphics[scale=0.25]{immElab.jpg}
        \subcaption{Immagine equalizzata}\label{esempio3}
    \end{minipage}
\end{figure}

\vspace{0.3cm}

È evidente come l'immagine \small (c) \normalsize presenti un contrasto maggiore rispetto alla \small(b)\normalsize. 
Questo è dato dall'ampliamento del range di valori assunti dai pixel dell'immagine, come evidenziato dalla tabella di memoria:
\begin{itemize}
    \item Immagine sorgente, addr. 2 - 17: i pixel assumono valori da 52 a 122;
    \item Immagine equalizzata, addr. 18 - 33: i pixel assumono valori da 0 a 255.
\end{itemize}

Non avendo ancora affrontato appieno il processo svolto dal componente, si ritiene opportuno non approfondire in questa sezione
l'insieme di passaggi che permettono la trasformazione dell'immagine riportata nell'esempio, che verranno ripresi in seguito. \par
Si rende disponibile un TestBench che replica l'elaborazione dell'immagine \small (b) \normalsize a questo \href{festadelsalame.com}{link}. \par

\vspace{2,5cm}

\section{Design del componente} %2
Si è scelto di descrivere un modulo \emph{single-process} tramite architettura \emph{behavioral} (comportamentale) in linguaggio VHDL.
Questo ha determinato la necessità di sviluppare un algoritmo adeguato allo svolgimento dell’operazione richiesta al componente, 
che può essere schematizzato secondo i seguenti passaggi chiave:

\begin{itemize}
    \item   [1.]    Lettura \texttt{n\_col} e \texttt{n\_row};
    \item   [2.]    Calcolo dimensione dell'immagine;
    \item   [3.]    Ciclo sui pixel dell’immagine sorgente per individuare tra di essi i valori di massimo e minimo;
    \item   [4.]    Calcolo dei valori necessari per l’elaborazione dell’immagine;
    \item   [5.]    Ciclo sui pixel dell’immagine sorgente per calcolare e salvare in memoria gli effettivi valori, pixel per pixel, dell’immagine equalizzata.
\end{itemize}

Il modulo prodotto opera quindi su una macchina a stati finiti che si occupa nel dettaglio dell’attuazione dell’algoritmo sviluppato.
\vspace{0,3cm}

\subsection{Macchina a stati finiti} %2.1

L’FSM da noi prodotta è composta da 10 stati, suddivisibili in 4 gruppi principali descritti di seguito.

\subsubsection{Stati ausiliari} %2.1.1
Gruppo di stati che realizza: inizio, richiesta di lettura, attesa della memoria e fine del processo.\par

\begin{itemize}
    \item [i.]      \textbf{\texttt{WT\_RST} - wait reset}: stato di attesa del segnale \texttt{i\_rst};
    \item [ii.]     \textbf{\texttt{WT\_STR} - wait start}: stato di attesa del segnale di \texttt{i\_start}. \\
                    In qualsiasi momento dell’elaborazione, se il segnale \texttt{i\_rst} è rilevato alto\footnotemark, anche non in corrispondenza di \texttt{i\_clk}, la macchina viene riportata in questo stato, tornando in attesa di un nuovo segnale di inizio elaborazione.\par
                    Al verificarsi della condizione \texttt{i\_start $=$ 1} vengono inizializzati tutti i valori necessari al processo, prima di passare allo stato successivo. Di particolare importanza per l’algoritmo sviluppato è il segnale \texttt{count} (inizialmente 0), che indica l’indirizzo a cui sarà effettuata l’operazione di read alla successiva richiesta di lettura del componente;
                    \footnotetext{In caso di reset si è supposto che il segnale \texttt{i\_start} venga riportato basso per il periodo in cui il \texttt{i\_rst} è alto.}
    \item [iii.]    \textbf{\texttt{RD\_REQ} - read request}: stato di abilitazione della memoria in lettura. Viene predisposto su \texttt{o\_address} l’indirizzo della RAM che deve essere letto;
    \item [iv.]     \textbf{\texttt{WT\_MEM} - wait memory}: stato di attesa della memoria che permette al valore richiesto in \texttt{RD\_REQ} di essere correttamente riportato sul segnale \texttt{i\_data} al ciclo di clock successivo. \par
                    È un nodo decisivo per l’FSM: viene costantemente rivisitato nei cicli di lettura dei pixel dell’immagine ed è responsabile del corretto instradamento del processo, grazie a condizioni su \texttt{count} e \texttt{shift\_value}. Si occupa inoltre dell’aggiornamento della variabile count stessa, e quindi della corretta gestione del successivo dato letto in memoria;
    \item [v.]      \textbf{\texttt{DONE} - done}: stato in cui \texttt{o\_done} viene posto a \texttt{‘1’} per segnalare la fine dell’elaborazione. A questo punto, si attende un valore di \texttt{i\_start} basso per tornare in \texttt{WT\_STR} e poter ricominciare il processo di equalizzazione di una nuova immagine;
\end{itemize}

\subsubsection{Calcolo dimensioni dell'immagine} %2.1.2
Gruppo di stati che permette il calcolo della dimensione effettiva dell’immagine da elaborare.\pagebreak

\begin{itemize}
    \item [vi.]     \textbf{\texttt{RD\_COL} - read column}: stato in cui il valore \texttt{n\_col} relativo all’immagine, pronto su \texttt{i\_data}, è salvato su una variabile temporanea per essere utilizzato in seguito;
    \item [vii.]    \textbf{\texttt{RD\_ROW} - read row}: stato in cui \texttt{n\_row}, pronto su \texttt{i\_data}, viene moltiplicato con il valore \texttt{n\_col} salvato precedentemente, per calcolare la dimensione effettiva dell’immagine e determinare se essa è adatta (\texttt{n\_col $\cdot$ n\_row $>$ 0}) o meno per il proseguimento dell’esecuzione;
\end{itemize}

\subsubsection{Ricerca dei valori di massimo e minimo dell'immagine} %2.1.3
Gruppo di stati che permette di individuare i valori minimo e massimo (\texttt{min} e \texttt{max} nel codice), necessari per l’effettiva elaborazione dell’immagine, tra quelli dei pixel dell’immagine sorgente.\par
La ricerca è svolta tramite un ciclo sui nodi \texttt{RD\_REQ}, \texttt{WT\_MEM} e \texttt{CMP\_DT} dell’FSM. Come già menzionato precedentemente, è \texttt{WT\_MEM} a occuparsi del corretto aggiornamento della variabile \texttt{count}, e quindi della lettura sequenziale dei pixel durante il ciclo.\par

\begin{itemize}
    \item [viii.]   \textbf{\texttt{CMP\_DT} - compare data}: stato in cui il valore del pixel dell’immagine relativo all’iterazione corrente, pronto sul segnale \texttt{i\_data}, viene confrontato con le variabili contenenti il minimo e massimo stabiliti fino a questa iterazione del ciclo di ricerca, aggiornandole se necessario.\par
                    Se la condizione di termine della ricerca (\texttt{count $\leq$ n\_col $\cdot$ n\_row $+$ 2}) si verifica, si riporta \texttt{count} pari all’indirizzo del primo pixel dell’immagine sorgente, ovvero \texttt{2}, per poi procedere alla fase di effettiva equalizzazione. Se invece non si è ancora scandagliata l’intera immagine in memoria, si procede nel ciclo al pixel successivo;
\end{itemize}

\subsubsection{Elaborazione dell'immagine} %2.1.4
Gruppo di stati che svolge l’effettiva elaborazione, pixel per pixel, dell’immagine da trasformare, tramite specifici valori calcolati in \texttt{PREP\_EL} e un ciclo sui nodi \texttt{RD\_REQ}, \texttt{WT\_MEM} e \texttt{EL\_DATA} della macchina a stati finiti.\par

\begin{itemize}
    \item   [ix.]   \textbf{\texttt{PREP\_EL} - prepare elaboration}: stato in cui vengono stabiliti, per mezzo dei dati ottenuti negli stati precedenti, il \texttt{delta\_value} e lo \texttt{shift\_level} relativi all’immagine da elaborare, necessari per il proseguimento del processo;
    \item   [x.]    \textbf{\texttt{EL\_DATA} - elaborate data}: stato in cui si svolge l’elaborazione effettiva del pixel dell’immagine relativo all’iterazione corrente. In particolare:
                    \begin{itemize}
                        \item [a.]  Si abilita in scrittura la memoria, ponendo in \texttt{o\_address} il valore dell’indirizzo di destinazione per la scrittura: \texttt{count $-$ 1 $+$ n\_col $\cdot$ n\_row}\footnotemark;
                                    \footnotetext{Essendo \texttt{count} già stato incrementato nello scorso passaggio dallo stato \texttt{WT\_MEM}, il pixel elaborato ad ogni iterazione è contenuto nell'indirizzo \texttt{count $-$ 1} della memoria.}
                        \item [b.]  Il valore del pixel dell’immagine originale, disponibile su \texttt{i\_data}, è utilizzato per calcolare quello del rispettivo pixel nell’immagine trasformata, il quale viene posto in \texttt{o\_data} per essere scritto in memoria;
                        \item [c.]  Se la condizione di termine dell’elaborazione dell’immagine (\texttt{count $\leq$ n\_col $\cdot$ n\_row $+$ 2}) si verifica, si passa a \texttt{DONE}, altrimenti si procede con l’equalizzazione del pixel successivo.
                    \end{itemize}
\end{itemize}

\pagebreak

%immagine FSM
\begin{figure}[ht] 
    \centering 
    \begin{tikzpicture}[->,>=stealth',shorten >=1pt,auto,node distance=3cm,
        semithick, initial text=$ $, initial where= above]

        \node[state, initial]           (0) {\texttt{WT\_RST}};
        \node[state, below of=0]        (1) {\texttt{WT\_STR}};
        \node[state, right=5cm of 1]    (2) {\texttt{RD\_COL}};
        \node[state, below of=1]        (3) {\texttt{RD\_ROW}};
        \node[state, below of=3]        (4) {\texttt{CPM\_DT}};
        \node[state, right=5cm of 3]    (5) {\texttt{RD\_REQ}};
        \node[state, right=5cm of 4]    (6) {\texttt{WT\_MEM}};
        \node[state, below of=4]        (7) {\texttt{PREP\_EL}};
        \node[state, below of=6]        (8) {\texttt{EL\_DATA}};
        \node[state, left of=7]         (9) {\texttt{DONE}};

        \draw   (0) edge                    node {\texttt{i\_rst=1}}                                            (1)
                (1) edge[loop left]         node {\texttt{i\_start=0}}                                          (1)
                (1) edge                    node [yshift=15pt, xshift=-38pt]{\texttt{i\_start=1}}               (5)
                (2) edge                    node {}                                                             (5)
                (3) edge                    node [xshift=-10pt] {\texttt{\emph{cond5}}}                         (5)
                (3) edge[left]              node [yshift=40pt, xshift=20pt]{\texttt{\emph{!cond5}}}             (9)
                (4) edge                    node [at start, yshift=8pt, xshift=25pt] {\texttt{\emph{cond6}}}    (5)
                (4) edge                    node[left] {\texttt{\emph{!cond6}}}                                 (7)
                (5) edge                    node {}                                                             (6)
                (6) edge[bend right, right] node {\texttt{\emph{cond1}}}                                        (2)
                (6) edge[right]             node [at start, xshift=-43pt] {\texttt{\emph{cond2}}}               (3)
                (6) edge                    node [at start, xshift=-15pt] {\texttt{\emph{cond3}}}               (4)
                (6) edge[left]              node {\texttt{\emph{cond4}}}                                        (8)                
                (7) edge                    node {}                                                             (5)
                (8) edge[bend right]        node[right]  {{\texttt{\emph{cond6}}}}                              (5)
                (8) edge[bend left=25]      node {{\texttt{\emph{!cond6}}}}                                     (9)  
                (9) edge[bend left]         node {\texttt{i\_start=0}}                                          (1)
                (9) edge[loop left]         node {\texttt{i\_start=1}}                                          (9);

    \end{tikzpicture}
    \caption{Diagramma degli stati della macchina a stati finiti utilizzata}
    \label{fig:fsm}
\end{figure}

\vspace{0,3cm}
Le condizioni utilizzate nel processo sono le seguenti:

\begin{center}
    \begin{tabular}{||c|c||}
        %\hline
        \hline
        \texttt{\emph{cond1}} & \texttt{count $=$ 0} \\
        \hline
        \texttt{\emph{cond2}} & \texttt{count $=$ 1} \\
        \hline
        \texttt{\emph{cond3}} & \texttt{shift\_level $=$ 3 (non ancora calcolato)} \\
        \hline
        \texttt{\emph{cond4}} & \texttt{\emph{!cond1} \&\& \emph{!cond2} \&\& \emph{!cond3}} \\   
        \hline
        \texttt{\emph{cond5}} & \texttt{n\_col $\cdot$ n\_row $>$ 0} \\
        \hline
        \texttt{\emph{cond6}} & \texttt{count $\leq$ n\_col $\cdot$ n\_row $+$ 2} \\      
        \hline
        %\hline
        \end{tabular}
\end{center}

\vspace{0,3cm}
Si ricorda inoltre, che per ogni stato dell'FSM è presente un arco uscente implicito diretto verso lo stato \texttt{WT\_STR}, 
che simboleggia la possibilità di interrompere in qualsiasi momento l'elaborazione dell'immagine corrente, tramite un segnale \texttt{i\_rst $=$ 1}.

\subsection{Approfondimento sull'equalizzazione dell'immagine} %2.2
\label{sec:appr}
La manipolazione del contrasto dell’immagine si fonda su 4 espressioni fondamentali. Si noti che le prime due sono valutate una sola volta per ogni immagine nello stato \texttt{PREP\_EL}, mentre le restanti sono determinate ad ogni iterazione del ciclo di elaborazione nello stato \texttt{EL\_DATA}.

\begin{itemize}
    \item   \texttt{delta\_value $=$ max $-$ min}\\
            \texttt{delta\_value} rappresenta la differenza tra il pixel più chiaro (valore maggiore – \texttt{max}) e il più scuro (valore minore – \texttt{min}) dell’immagine;
    \item   \texttt{shift\_value $=$ (8 $-$ floor\footnotemark($\log_{2} $(delta\_value $+$ 1)))}\\
            \texttt{shift\_value} determina il numero di shift a sinistra da applicare al risultato della differenza tra il pixel considerato nell’iterazione corrente e il pixel di valore minore dell’immagine;
            \footnotetext{La funzione \texttt{floor(x)} svolge l’arrotondamento per difetto del valore x fornitogli come argomento.}

            \vspace{0,3cm}
            \begin{table}[h]
                \centering
                \def\arraystretch{1.3} %un po' di padding
                \caption{Esempi di valori di \texttt{delta\_value}, \texttt{floor(x)} e \texttt{shift\_level} possibili per la codifica delle immagini in scala di grigi a 256 livelli.}
                \begin{tabular}{||c|c|c||}
                    \hline
                    \texttt{delta\_value} & \texttt{floor(x)\footnotemark} & \texttt{shift\_level}\\
                    \hline \hline
                    0       & 0         & 8         \\\hline
                    1       & 1         & 7         \\\hline
                    2       & 1         & 7         \\\hline
                    3       & 2         & 6         \\\hline
                    \vdots  & \vdots    & \vdots    \\\hline
                    6       & 2         & 6         \\\hline
                    7       & 3         & 5         \\\hline
                    8       & 3         & 5         \\\hline
                    \vdots  & \vdots    & \vdots    \\\hline
                    14      & 3         & 5         \\\hline
                    15      & 4         & 4         \\\hline
                    16      & 4         & 4         \\\hline
                    \vdots  & \vdots    & \vdots    \\\hline
                    29      & 4         & 4         \\\hline
                \end{tabular}\hspace{20pt}
                \begin{tabular}{||c|c|c||}
                    \hline
                    \texttt{delta\_value} & \texttt{floor(x)\footnotemark[\value{footnote}]} & \texttt{shift\_level}\\
                    \hline \hline
                    30      & 4         & 4         \\\hline
                    31      & 5         & 3         \\\hline
                    32      & 5         & 3         \\\hline
                    \vdots  & \vdots    & \vdots    \\\hline
                    62      & 5         & 3         \\\hline
                    63      & 6         & 2         \\\hline
                    64      & 6         & 2         \\\hline
                    \vdots  & \vdots    & \vdots    \\\hline
                    126     & 6         & 2         \\\hline
                    127     & 7         & 1         \\\hline
                    128     & 7         & 1         \\\hline
                    \vdots  & \vdots    & \vdots    \\\hline
                    254     & 7         & 1         \\\hline
                    255     & 8         & 0         \\\hline
                \end{tabular}
            \end{table}
            \footnotetext{Dove: \texttt{x $=$ $\log_{2} $(delta\_value $+$ 1)}.}
            \vspace{0,3cm}

    \item   \texttt{temp\_pixel $=$ (current\_pixel $-$ min) $<<$ shift\_level}\\
            \texttt{temp\_pixel} rappresenta il possibile valore da attribuire nell’immagine finale al pixel valutato in questa iterazione del ciclo di elaborazione. Si noti che il suo valore potrebbe superare il limite massimo di 255 imposto dalla codifica in scala di grigi a 256 livelli. Per questo motivo non può essere utilizzato \emph{“as-is”}.
    \item   \texttt{new\_pixel $=$ min(255, temp\_pixel)}\\
            \texttt{new\_pixel} è pari al minimo tra \texttt{255} e \texttt{temp\_pixel} e rappresenta l’effettivo valore che verrà scritto sulla RAM per il pixel considerato in questa iterazione.
\end{itemize}\par

Riprendendo quindi l'esempio proposto nella \hyperref[sec:esempio]{Sezione 1.4}:
\begin{itemize}
    \item   \texttt{delta\_value $=$ max $-$ min = 122 - $-$ 52 = 70}\\
    \item   \texttt{shift\_value $=$ (8 $-$ floor($\log_{2} $(delta\_value $+$ 1))) $=$ 8 - floor(6.149) = 2}\\
\end{itemize}

Iterando quindi sui pixel dell'immagine sorgente:\\

\begin{center}
    \def\arraystretch{1.3} %un po' di padding
    \begin{tabular}{||c|c|c||} 
        \hline
        \texttt{current\_pixel} & \texttt{temp\_pixel} & \texttt{new\_pixel}\\
        \hline \hline
        52       & 0         & 0         \\\hline
        55       & 12        & 12        \\\hline
        122      & 5         & 255       \\\hline
        70       & 72        & 72        \\\hline
        \vdots   & \vdots    & \vdots    \\\hline
    \end{tabular}        
\end{center}

\vspace{0,3cm}

\subsection{Scelte progettuali e ottimizzazioni} %2.3
Si è scelto di progettare un componente sensibile al clock su \emph{rising\_edge}.\par
Nell’implementazione dell’algoritmo si sottolinea la scelta di mantenere l’operazione di moltiplicazione per il calcolo della dimensione effettiva dell’immagine da processare. Nonostante il prodotto sia un operatore pesante rispetto alla semplice somma o differenza, lavorando su segnali a 8 bit la sintesi è capace di gestirlo in modo efficace tramite l’inserimento di alcuni \emph{DPS}\footnotemark.
Si noti inoltre che l’operazione di moltiplicazione ricorre una sola volta nell’intera elaborazione di ogni singola immagine, e non risulta quindi particolarmente rilevante rispetto all’intero processo. Questo è stato verificato attraverso lo sviluppo simultaneo rispetto alla versione proposta di una ulteriore implementazione che non facesse uso di nessun prodotto all’interno del processo. In quest’ultima non si sono osservati vantaggi significativi in termini di tempo e area di sintesi, a scapito della leggibilità del codice stesso.\par
Da ultimo, si noti che \texttt{shift\_level}, assumendo valore intero compreso tra 0 e 8, è facilmente ricavabile tramite controlli di soglia. Questo permette di evitare nel codice l’effettivo sviluppo del logaritmo indicato nella realtiva formula matematica, enunciata alla \hyperref[sec:appr]{Sezione 2.2}.\par 
Di seguito si riporta il blocco IF/ELSE che si è utilizzato per il calcolo dello \texttt{shift\_value} relativo a un dato \texttt{delta\_value}:
\footnotetext{DSP (digital signal processor) è un processore dedicato e ottimizzato per eseguire in maniera estremamente efficiente sequenze di istruzioni ricorrenti (come ad esempio somme, moltiplicazioni e traslazioni) nell'elaborazione di segnali digitali.}

\begin{minted}{VHDL}
    temp_integer := TO_INTEGER(unsigned (max)) - TO_INTEGER(unsigned (min)); --delta_value
    IF temp_integer = 0 THEN
        shift_level <= 8;
    ELSIF temp_integer > 0 AND temp_integer < 3 THEN
        shift_level <= 7;
    ELSIF temp_integer > 2 AND temp_integer < 7 THEN
        shift_level <= 6;
    ELSIF temp_integer > 6 AND temp_integer < 15 THEN
        shift_level <= 5;
    ELSIF temp_integer > 14 AND temp_integer < 31 THEN
        shift_level <= 4;
    ELSIF temp_integer > 30 AND temp_integer < 63 THEN
        shift_level <= 3;
    ELSIF temp_integer > 62 AND temp_integer < 127 THEN
        shift_level <= 2;
    ELSIF temp_integer > 126 AND temp_integer < 255 THEN
        shift_level <= 1;
    ELSE
        shift_level <= 0;
    END IF;
\end{minted}

\vspace{2.5cm}

\end{document}
